\section*{Anexo: Corpus completo de 42 bibliotecas}

\begin{longtable}{rlrll}
\caption{Corpus completo: 42 bibliotecas hispánicas (1518-1811)}\\
\toprule
\textbf{Nº} & \textbf{Propietario} & \textbf{Obras} & \textbf{Fecha} & \textbf{Perfil social} \\
\midrule
\endfirsthead
\multicolumn{5}{c}{\tablename\ \thetable\ -- Continuación}\\
\toprule
\textbf{Nº} & \textbf{Propietario} & \textbf{Obras} & \textbf{Fecha} & \textbf{Perfil social} \\
\midrule
\endhead
\midrule
\multicolumn{5}{r}{\textit{Continúa en la página siguiente...}}\\
\endfoot
\bottomrule
\textbf{TOTAL: 42 bibliotecas, 54.059 obras catalogadas}\\
\endlastfoot

\rowcolor{yellow!20}
1 & Ramírez de Prado.  Lorenzo & 8.951 & 1662 & Magistrado (Consejo) \\
\rowcolor{yellow!20}
2 & Velasco y Ceballos. Fernando José de & 7.323 & 1791 & Camarista de Castilla \\
\rowcolor{yellow!20}
3 & Monasterio de San Martín & 7.119 & 1788 & Institución eclesiástica \\
\rowcolor{yellow!20}
4 & Sarmiento Acuña.  Diego. conde de Gondomar & 6.471 & +1626 & Noble + Diplomático \\
\rowcolor{yellow!20}
5 & Rodríguez de Campomanes. Pedro & 5.500 & †1802 & Fiscal del Consejo \\
\rowcolor{yellow!20}
6 & Jovellanos. Gaspar Melchor de & 5.374 & †1811 & Ministro ilustrado \\
7 & Biblioteca de la Torre Alta del Alcázar & 2.012 & 1637 & Biblioteca real \\
8 & Cerda. Antonio Juan Luis de la. VII duque de Medinaceli & 1.474 & 1673 & Alta nobleza (Duque) \\
9 & Lastanosa. Vincencio Juan de & 974 & 1662 & Erudito + Noble \\
10 & Aragón. Fernando de. duque de Calabria & 793 & 1550 & Alta nobleza (Duque) \\
11 & Frago. Pedro del. obispo de Cuenca & 740 & 1584 & Alto clero (Obispo) \\
12 & Osorio. Alonso. VII marqués de Astorga. Inventario A & 707 & 1573 & Alta nobleza (Marqués) \\
13 & Mendoza. Mencía de & 616 & 1553, 1562 & Alta nobleza femenina \\
14 & Díaz de Vivar y Mendoza. Rodrigo. I Marqués del Cenete & 582 & 1523 & Alta nobleza (Marqués) \\
15 & Caro. Rodrigo & 530 & 1647 & Erudito + Arqueólogo \\
16 & Cromberger. Juan. impresor & 525 & 1540 & Impresor profesional \\
17 & Garza Falcón. Joseph Manuel de la. oidor de la Audiencia de la Nueva Galicia & 514 & 1763 & Magistrado colonial \\
18 & Barrientos. Bartolomé. maestro de artes liberales & 492 & 1576 & Maestro artes liberales \\
19 & Osorio. Alonso. VII marqués de Astorga. Inventario B & 354 & 1589 & Alta nobleza (Marqués) \\
20 & Simón Abril. Pedro & 333 & +1595 & Humanista + Pedagogo \\
21 & Fernández de Córdoba. Pedro. marqués de Priego & 268 & 1518 & Alta nobleza (Marqués) \\
22 & Enríquez de Ribera. Fadrique. I marqués de Tarifa & 221 & 1532 & Alta nobleza (Marqués) \\
23 & Villalba. Felipa de & 216 & 1626 & Nobleza femenina \\
24 & Quevedo. Francisco de & 189 & 1647 & Escritor + Noble \\
25 & Inca Garcilaso de la Vega & 188 & +1616 & Escritor + Cronista \\
26 & Arias Dávila y Bobadilla. Francisco. IV conde de Puñonrostro & 182 & 1610 & Nobleza (Conde) \\
27 & Argote de Molina. Gonzalo & 166 & +1596 & Erudito + Historiador \\
28 & Velázquez. Diego & 154 & 1660 & Pintor de cámara \\
29 & López de Fuentesdaño. Juan. inquisidor de Valladolid & 152 & 1631 & Inquisidor \\
30 & Barros. Alonso de & 151 & +1604 & Militar + Escritor \\
31 & Villasinda. Juan Alonso de. relator & 141 & 1631 & Relator real \\
32 & Gómez de Silva. Ruy. III duque de Pastrana & 94 & 1626 & Alta nobleza (Duque) \\
33 & Mendoza. Bernardino de & 79 & +1604 & Militar + Diplomático \\
34 & Schomburg. Juan Carlos. conde de Schomburg & 79 & 1640 & Nobleza extranjera \\
35 & Brizuela. Mateo de & 74 & -1557 & Nobleza \\
36 & Rojas. Antonio de. caballero cristiano & 64 & 1556 & Caballero \\
37 & Hurtado de Mendoza. Diego. I conde de Mélito & 60 & 1536 & Alta nobleza (Conde) \\
38 & Castro Enríquez Osorio. Beatriz de. condesa de Lemos & 48 & 1570 & Alta nobleza (Marqués) \\
39 & Theotokópoulos. Doménikos. el Greco & 47 & 1614 & Pintor (El Greco) \\
40 & Gómez de Silva. Ruy. príncipe de Éboli y I duque de Pastrana & 40 & 1573 & Alta nobleza (Príncipe) \\
41 & Rodríguez de la Torre. Francisco. secretario del rey Carlos II & 31 & 1699 & Secretario real \\
42 & Silva y Mendoza. Rodrigo de. II duque de Pastrana & 31 & 1596 & Alta nobleza (Duque) \\
\end{longtable}

\textbf{Nota metodológica:} Las 6 megabibliotecas (>5.000 obras, resaltadas en amarillo) concentran 40.937 obras (75,8\% del total), evidenciando una extrema concentración del capital cultural librario en la cúspide de la élite magistral e ilustrada.
